\documentclass[12pt]{article}

% Idioma e codificação
\usepackage[brazil]{babel}
\usepackage[utf8]{inputenc}
\usepackage[T1]{fontenc}

% Matemática e listas
\usepackage{amsmath, amssymb}
\usepackage{enumitem}
\usepackage{float}
\usepackage{colortbl}
\usepackage{geometry}
\usepackage{graphicx}
\usepackage{xcolor}

% Cores institucionais e personalizadas
\definecolor{maincolor}{RGB}{0,100,50}
\definecolor{forestgreen}{RGB}{0,100,50}
\definecolor{bgsql}{RGB}{245, 245, 245}
\definecolor{verdeuni}{RGB}{0, 100, 0}
\definecolor{blueuni}{RGB}{0, 0, 128}

% Cabeçalho e rodapé
\usepackage{fancyhdr}
\usepackage{tabularx}
\usepackage{booktabs}
\usepackage[most]{tcolorbox}
\tcbuselibrary{minted}

\newtcblisting{sqlcode}{
    listing engine=minted,
    minted language=sql,
    colback=bgsql,
    colframe=verdeuni,
    listing only,
    title=Gabarito SQL,
    fonttitle=\bfseries\color{white},
    coltitle=white, 
    colbacktitle=verdeuni
}

\geometry{margin=2.5cm}
\pagestyle{fancy}
\fancyhf{}
\renewcommand{\headrulewidth}{2pt}
\renewcommand{\footrulewidth}{2pt}
\renewcommand{\headrule}{\hbox to\headwidth{\color{forestgreen}\leaders\hrule height \headrulewidth\hfill}}
\renewcommand{\footrule}{\hbox to\headwidth{\color{forestgreen}\leaders\hrule height \footrulewidth\hfill}}
\setlength{\headheight}{52pt}

\fancyhead[L]{\textbf{UniRV}}
\fancyhead[R]{%
  \raggedleft
  \textcolor{forestgreen}{\small \textbf{Gabarito - Lista Complementar}\\Me.~Sergio Souza Novak\\\textit{Banco de Dados}}
}
\fancyfoot[C]{\textcolor{forestgreen}{\thepage}}

\begin{document}

\section*{Gabarito Oficial – Criação, Manipulação e Análise}

\hrulefill

\section*{Exercício 1 – Gestão de Estoque}



\begin{sqlcode}
-- 1. Criação da Tabela
CREATE TABLE produto (
    id SERIAL PRIMARY KEY,
    nome VARCHAR(100) NOT NULL,
    descricao TEXT,
    preco DECIMAL(10,2) NOT NULL,
    estoque INTEGER DEFAULT 0,
    sku CHAR(8) UNIQUE NOT NULL,
    criado_em TIMESTAMP DEFAULT CURRENT_TIMESTAMP
);

-- 2. Inserção de Dados
INSERT INTO produto (nome, preco, estoque, sku) VALUES
('Teclado Mecânico', 250.00, 15, 'TEC12345'),
('Mouse Gamer', 150.00, 5, 'MOU55443'),
('Monitor 24pol', 890.00, 0, 'MON88776');

-- 3. Consulta DQL
SELECT nome, preco 
FROM produto 
WHERE estoque > 0 
ORDER BY preco DESC;
\end{sqlcode}

\hrulefill

\section*{Exercício 2 – Relatório de Vendas}



\begin{sqlcode}
-- 1. Quantidade de vendas por vendedor
SELECT vendedor_id, COUNT(id) AS total_vendas
FROM venda
GROUP BY vendedor_id;

-- 2. Filtro com HAVING
SELECT vendedor_id, COUNT(id) AS total_vendas
FROM venda
GROUP BY vendedor_id
HAVING COUNT(id) > 10;
\end{sqlcode}

\hrulefill

\section*{Exercício 3 – Status de Clientes}



\begin{sqlcode}
-- 1. Criação e Inserção
CREATE TABLE cliente (
    id SERIAL PRIMARY KEY,
    nome VARCHAR(100) NOT NULL,
    cpf CHAR(11) UNIQUE NOT NULL,
    email VARCHAR(150) UNIQUE,
    ativo BOOLEAN DEFAULT TRUE
);

INSERT INTO cliente (nome, cpf, ativo) VALUES 
('Ana', '11122233344', true),
('Bob', '55566677788', true),
('Eva', '99900011122', false);

-- 2. Contagem por status
SELECT ativo, COUNT(id) 
FROM cliente 
GROUP BY ativo;

-- 3. Filtro de grupo (menos de 2 clientes)
SELECT ativo, COUNT(id) 
FROM cliente 
GROUP BY ativo
HAVING COUNT(id) < 2;
\end{sqlcode}
\textbf{Justificativa:} Usamos \texttt{HAVING} porque o filtro é aplicado sobre o resultado de uma função agregada (\texttt{COUNT}), o que ocorre após o agrupamento dos dados.

\hrulefill

\section*{Exercício 4 – Filtros Compostos em Frota}



\begin{sqlcode}
-- 1. Contagem específica com WHERE
SELECT COUNT(id) 
FROM carro 
WHERE marca = 'Toyota' AND ano > 2020;

-- 2. Agrupamento e Ordenação
SELECT marca, COUNT(id) AS total
FROM carro
GROUP BY marca
ORDER BY marca ASC;
\end{sqlcode}

\hrulefill

\section*{Exercício 5 – Folha de Pagamento}



\begin{sqlcode}
-- 1. Soma total de funcionários ativos
SELECT SUM(salario) AS custo_total_folha
FROM empregado
WHERE ativo = TRUE;
\end{sqlcode}
\textbf{Resposta 2:} Usamos \texttt{SUM} para somar os valores monetários contidos na coluna salário. O \texttt{COUNT} apenas contaria quantas linhas existem, ignorando o valor do salário de cada funcionário.

\hrulefill

\section*{Desafio de Lógica – Correção de Sintaxe}



\begin{sqlcode}
-- Código Corrigido
SELECT departamento, COUNT(id)
FROM funcionarios
WHERE salario > 2000
GROUP BY departamento
HAVING COUNT(id) > 5;
\end{sqlcode}

\textbf{Explicação:} O \texttt{WHERE} deve vir antes do \texttt{GROUP BY} porque ele atua como um filtro de linhas individuais (reduzindo o volume de dados a ser processado). Somente após esse filtro o banco de dados organiza o que sobrou em grupos para então aplicar o \texttt{HAVING}.

\end{document}